% Options for packages loaded elsewhere
\PassOptionsToPackage{unicode}{hyperref}
\PassOptionsToPackage{hyphens}{url}
%
\documentclass[
]{article}
\usepackage{amsmath,amssymb}
\usepackage{iftex}
\ifPDFTeX
  \usepackage[T1]{fontenc}
  \usepackage[utf8]{inputenc}
  \usepackage{textcomp} % provide euro and other symbols
\else % if luatex or xetex
  \usepackage{unicode-math} % this also loads fontspec
  \defaultfontfeatures{Scale=MatchLowercase}
  \defaultfontfeatures[\rmfamily]{Ligatures=TeX,Scale=1}
\fi
\usepackage{lmodern}
\ifPDFTeX\else
  % xetex/luatex font selection
\fi
% Use upquote if available, for straight quotes in verbatim environments
\IfFileExists{upquote.sty}{\usepackage{upquote}}{}
\IfFileExists{microtype.sty}{% use microtype if available
  \usepackage[]{microtype}
  \UseMicrotypeSet[protrusion]{basicmath} % disable protrusion for tt fonts
}{}
\makeatletter
\@ifundefined{KOMAClassName}{% if non-KOMA class
  \IfFileExists{parskip.sty}{%
    \usepackage{parskip}
  }{% else
    \setlength{\parindent}{0pt}
    \setlength{\parskip}{6pt plus 2pt minus 1pt}}
}{% if KOMA class
  \KOMAoptions{parskip=half}}
\makeatother
\usepackage{xcolor}
\usepackage[margin=1in]{geometry}
\usepackage{color}
\usepackage{fancyvrb}
\newcommand{\VerbBar}{|}
\newcommand{\VERB}{\Verb[commandchars=\\\{\}]}
\DefineVerbatimEnvironment{Highlighting}{Verbatim}{commandchars=\\\{\}}
% Add ',fontsize=\small' for more characters per line
\usepackage{framed}
\definecolor{shadecolor}{RGB}{248,248,248}
\newenvironment{Shaded}{\begin{snugshade}}{\end{snugshade}}
\newcommand{\AlertTok}[1]{\textcolor[rgb]{0.94,0.16,0.16}{#1}}
\newcommand{\AnnotationTok}[1]{\textcolor[rgb]{0.56,0.35,0.01}{\textbf{\textit{#1}}}}
\newcommand{\AttributeTok}[1]{\textcolor[rgb]{0.13,0.29,0.53}{#1}}
\newcommand{\BaseNTok}[1]{\textcolor[rgb]{0.00,0.00,0.81}{#1}}
\newcommand{\BuiltInTok}[1]{#1}
\newcommand{\CharTok}[1]{\textcolor[rgb]{0.31,0.60,0.02}{#1}}
\newcommand{\CommentTok}[1]{\textcolor[rgb]{0.56,0.35,0.01}{\textit{#1}}}
\newcommand{\CommentVarTok}[1]{\textcolor[rgb]{0.56,0.35,0.01}{\textbf{\textit{#1}}}}
\newcommand{\ConstantTok}[1]{\textcolor[rgb]{0.56,0.35,0.01}{#1}}
\newcommand{\ControlFlowTok}[1]{\textcolor[rgb]{0.13,0.29,0.53}{\textbf{#1}}}
\newcommand{\DataTypeTok}[1]{\textcolor[rgb]{0.13,0.29,0.53}{#1}}
\newcommand{\DecValTok}[1]{\textcolor[rgb]{0.00,0.00,0.81}{#1}}
\newcommand{\DocumentationTok}[1]{\textcolor[rgb]{0.56,0.35,0.01}{\textbf{\textit{#1}}}}
\newcommand{\ErrorTok}[1]{\textcolor[rgb]{0.64,0.00,0.00}{\textbf{#1}}}
\newcommand{\ExtensionTok}[1]{#1}
\newcommand{\FloatTok}[1]{\textcolor[rgb]{0.00,0.00,0.81}{#1}}
\newcommand{\FunctionTok}[1]{\textcolor[rgb]{0.13,0.29,0.53}{\textbf{#1}}}
\newcommand{\ImportTok}[1]{#1}
\newcommand{\InformationTok}[1]{\textcolor[rgb]{0.56,0.35,0.01}{\textbf{\textit{#1}}}}
\newcommand{\KeywordTok}[1]{\textcolor[rgb]{0.13,0.29,0.53}{\textbf{#1}}}
\newcommand{\NormalTok}[1]{#1}
\newcommand{\OperatorTok}[1]{\textcolor[rgb]{0.81,0.36,0.00}{\textbf{#1}}}
\newcommand{\OtherTok}[1]{\textcolor[rgb]{0.56,0.35,0.01}{#1}}
\newcommand{\PreprocessorTok}[1]{\textcolor[rgb]{0.56,0.35,0.01}{\textit{#1}}}
\newcommand{\RegionMarkerTok}[1]{#1}
\newcommand{\SpecialCharTok}[1]{\textcolor[rgb]{0.81,0.36,0.00}{\textbf{#1}}}
\newcommand{\SpecialStringTok}[1]{\textcolor[rgb]{0.31,0.60,0.02}{#1}}
\newcommand{\StringTok}[1]{\textcolor[rgb]{0.31,0.60,0.02}{#1}}
\newcommand{\VariableTok}[1]{\textcolor[rgb]{0.00,0.00,0.00}{#1}}
\newcommand{\VerbatimStringTok}[1]{\textcolor[rgb]{0.31,0.60,0.02}{#1}}
\newcommand{\WarningTok}[1]{\textcolor[rgb]{0.56,0.35,0.01}{\textbf{\textit{#1}}}}
\usepackage{graphicx}
\makeatletter
\def\maxwidth{\ifdim\Gin@nat@width>\linewidth\linewidth\else\Gin@nat@width\fi}
\def\maxheight{\ifdim\Gin@nat@height>\textheight\textheight\else\Gin@nat@height\fi}
\makeatother
% Scale images if necessary, so that they will not overflow the page
% margins by default, and it is still possible to overwrite the defaults
% using explicit options in \includegraphics[width, height, ...]{}
\setkeys{Gin}{width=\maxwidth,height=\maxheight,keepaspectratio}
% Set default figure placement to htbp
\makeatletter
\def\fps@figure{htbp}
\makeatother
\setlength{\emergencystretch}{3em} % prevent overfull lines
\providecommand{\tightlist}{%
  \setlength{\itemsep}{0pt}\setlength{\parskip}{0pt}}
\setcounter{secnumdepth}{-\maxdimen} % remove section numbering
\ifLuaTeX
  \usepackage{selnolig}  % disable illegal ligatures
\fi
\usepackage{bookmark}
\IfFileExists{xurl.sty}{\usepackage{xurl}}{} % add URL line breaks if available
\urlstyle{same}
\hypersetup{
  hidelinks,
  pdfcreator={LaTeX via pandoc}}

\author{}
\date{\vspace{-2.5em}}

\begin{document}

\begin{Shaded}
\begin{Highlighting}[]
\CommentTok{\# Aleatorización avanzada: Parámetros y variantes para dimensiones y contextos}

\CommentTok{\# Contextualización: nombres del objeto y descripción}
\NormalTok{objetos\_posibles }\OtherTok{\textless{}{-}} \FunctionTok{c}\NormalTok{(}\StringTok{"parabrisas"}\NormalTok{, }\StringTok{"ventanas laterales"}\NormalTok{, }\StringTok{"pantallas de instrumentos"}\NormalTok{)}
\NormalTok{objeto }\OtherTok{\textless{}{-}} \FunctionTok{sample}\NormalTok{(objetos\_posibles,}\DecValTok{1}\NormalTok{) }\CommentTok{\# Contexto principal}

\CommentTok{\# Dimensiones del parabrisas: dos medidas principales (largo x ancho)}
\NormalTok{largo\_parabrisas }\OtherTok{\textless{}{-}} \DecValTok{13} \CommentTok{\# cm fijo según enunciado de la imagen}
\CommentTok{\# Aleatorizamos el ancho principal, compatible con valores mostrados en imagen}
\NormalTok{ancho\_parabrisas }\OtherTok{\textless{}{-}} \FunctionTok{sample}\NormalTok{(}\FunctionTok{c}\NormalTok{(}\DecValTok{8}\NormalTok{, }\DecValTok{9}\NormalTok{, }\DecValTok{10}\NormalTok{), }\DecValTok{1}\NormalTok{) }\CommentTok{\# para variante}

\CommentTok{\# Plumillas (escobillas): número, longitud y apertura angular}
\NormalTok{num\_plumillas }\OtherTok{\textless{}{-}} \DecValTok{2} \CommentTok{\# fijo en problema}

\CommentTok{\# Aleatorizamos la longitud de las plumillas (en cm, rango plausible 5{-}7)}
\NormalTok{longitud\_plumilla }\OtherTok{\textless{}{-}} \FunctionTok{sample}\NormalTok{(}\FunctionTok{c}\NormalTok{(}\DecValTok{6}\NormalTok{,}\DecValTok{7}\NormalTok{), }\DecValTok{1}\NormalTok{) }

\CommentTok{\# Ángulo de apertura de las plumillas: fijo 90 grados en imagen,}
\CommentTok{\# se admite variación para aleatorización pero aquí fijo para precisión del problema original.}
\NormalTok{angulo\_apertura }\OtherTok{\textless{}{-}} \DecValTok{90}

\CommentTok{\# Configuración de opciones de dibujo y estilo para TikZ}
\NormalTok{color\_linea\_tabla }\OtherTok{\textless{}{-}} \StringTok{"red!80!black"}
\NormalTok{color\_plumilla }\OtherTok{\textless{}{-}} \StringTok{"black"}
\NormalTok{ancho\_linea }\OtherTok{\textless{}{-}} \StringTok{"thick"}

\CommentTok{\# Numeración de opciones y texto}
\CommentTok{\# Para aleatorización completa se puede permutar opciones, pero aquí igual para preservar opción correcta señalada.}

\CommentTok{\# Variables para las dimensiones presentadas en opciones (para validar en solución):}
\CommentTok{\# Opciones equivalentes a las de la imagen — a, b, c, d — con distinta correspondencia:}

\CommentTok{\# Opción A:}
\CommentTok{\# alto caja 13cm, ancho caja 8 cm; plumillas 6 cm + 6 cm}

\CommentTok{\# Opción B:}
\CommentTok{\# alto caja 8cm, ancho caja 13 cm; plumillas 6 cm y 8 cm}

\CommentTok{\# Opción C:}
\CommentTok{\# alto caja 8cm, ancho caja 13 cm; plumillas 6 cm y 6 cm}

\CommentTok{\# Opción D:}
\CommentTok{\# alto caja 6 cm, ancho caja 13 cm; plumillas 8 cm y 8 cm}

\CommentTok{\# Los datos para cada opción (más paramétricos y aleatorios):}
\NormalTok{opciones\_lista }\OtherTok{\textless{}{-}} \FunctionTok{list}\NormalTok{(}
  \AttributeTok{a =} \FunctionTok{list}\NormalTok{(}\AttributeTok{alto =} \DecValTok{13}\NormalTok{, }\AttributeTok{ancho =} \DecValTok{8}\NormalTok{, }\AttributeTok{plumilla1 =} \DecValTok{6}\NormalTok{, }\AttributeTok{plumilla2 =} \DecValTok{6}\NormalTok{),}
  \AttributeTok{b =} \FunctionTok{list}\NormalTok{(}\AttributeTok{alto =} \DecValTok{8}\NormalTok{, }\AttributeTok{ancho =} \DecValTok{13}\NormalTok{, }\AttributeTok{plumilla1 =} \DecValTok{6}\NormalTok{, }\AttributeTok{plumilla2 =} \DecValTok{8}\NormalTok{),}
  \AttributeTok{c =} \FunctionTok{list}\NormalTok{(}\AttributeTok{alto =} \DecValTok{8}\NormalTok{, }\AttributeTok{ancho =} \DecValTok{13}\NormalTok{, }\AttributeTok{plumilla1 =} \DecValTok{6}\NormalTok{, }\AttributeTok{plumilla2 =} \DecValTok{6}\NormalTok{),}
  \AttributeTok{d =} \FunctionTok{list}\NormalTok{(}\AttributeTok{alto =} \DecValTok{6}\NormalTok{, }\AttributeTok{ancho =} \DecValTok{13}\NormalTok{, }\AttributeTok{plumilla1 =} \DecValTok{8}\NormalTok{, }\AttributeTok{plumilla2 =} \DecValTok{8}\NormalTok{)}
\NormalTok{)}

\CommentTok{\# Correcta es la opción a:}
\NormalTok{opcion\_correcta }\OtherTok{\textless{}{-}} \StringTok{"a"}

\CommentTok{\# Variables para labels y títulos formateados}
\NormalTok{titulo\_parabrisas }\OtherTok{\textless{}{-}} \FunctionTok{paste0}\NormalTok{(}\StringTok{"Dimensiones del "}\NormalTok{, objeto)}
\end{Highlighting}
\end{Shaded}

\section{Question}\label{question}

Las dimensiones de un ventanas laterales son 13 cm de largo y 9 cm de
ancho.

Este objeto es plano y de forma rectangular, correspondiente a un
carrito de juguete que cuenta con 2 plumillas de 6 cm cada una para
limpiar esta superficie. Estas plumillas tienen un ángulo de apertura de
90°, tal como se muestra en la siguiente ilustración.

¿Cuál de las siguientes opciones representa correctamente el esquema que
modela las dimensiones del ventanas laterales y sus plumillas?

\subsection{Answerlist}\label{answerlist}

\begin{tikzpicture}[scale=0.6]
\draw[thick, color=red!80!black] (0,0) rectangle (8,13);
\draw[|-|, color=red!80!black] (-0.4,0) -- node[left]{13 cm} (-0.4,13);
\draw[|-|, color=red!80!black] (0,-0.4) -- node[below]{8 cm} (8,-0.4);
\draw[thick] (0,0) arc (180:90:6);
\draw[thick] (6,0) -- (0,0);
\draw[thick] (0,6) -- (0,0);
\draw[thick] (8,0) arc (0:90:6);
\draw[thick] (2,0) -- (8,0);
\draw[thick] (8,6) -- (8,0);
\node at (4, -1) {(A));
\end{tikzpicture}

\begin{tikzpicture}[scale=0.6]
\draw[thick, color=red!80!black] (0,0) rectangle (13,8);
\draw[|-|, color=red!80!black] (-0.4,0) -- node[left]{8 cm} (-0.4,8);
\draw[|-|, color=red!80!black] (0,-0.4) -- node[below]{13 cm} (13,-0.4);
\draw[thick] (0,0) arc (180:90:6);
\draw[thick] (6,0) -- (0,0);
\draw[thick] (0,6) -- (0,0);
\draw[thick] (13,0) arc (0:90:8);
\draw[thick] (5,0) -- (13,0);
\draw[thick] (13,8) -- (13,0);
\node at (6.5, -1) {(B));
\end{tikzpicture}

\begin{tikzpicture}[scale=0.6]
\draw[thick, color=red!80!black] (0,0) rectangle (13,8);
\draw[|-|, color=red!80!black] (-0.4,0) -- node[left]{8 cm} (-0.4,8);
\draw[|-|, color=red!80!black] (0,-0.4) -- node[below]{13 cm} (13,-0.4);
\draw[thick] (0,0) arc (180:90:6);
\draw[thick] (6,0) -- (0,0);
\draw[thick] (0,6) -- (0,0);
\draw[thick] (13,0) arc (0:90:6);
\draw[thick] (7,0) -- (13,0);
\draw[thick] (13,6) -- (13,0);
\node at (6.5, -1) {(C));
\end{tikzpicture}

\begin{tikzpicture}[scale=0.6]
\draw[thick, color=red!80!black] (0,0) rectangle (13,6);
\draw[|-|, color=red!80!black] (-0.4,0) -- node[left]{6 cm} (-0.4,6);
\draw[|-|, color=red!80!black] (0,-0.4) -- node[below]{13 cm} (13,-0.4);
\draw[thick] (0,0) arc (180:90:8);
\draw[thick] (8,0) -- (0,0);
\draw[thick] (0,8) -- (0,0);
\draw[thick] (13,0) arc (0:90:8);
\draw[thick] (5,0) -- (13,0);
\draw[thick] (13,8) -- (13,0);
\node at (6.5, -1) {(D));
\end{tikzpicture}

\section{Solution}\label{solution}

La opción correcta es la que representa las dimensiones originales dadas
y las plumillas tal y como se describen:

\begin{itemize}
\tightlist
\item
  Largo (alto del rectángulo) 13 cm
\item
  Ancho (ancho del rectángulo) 8 cm
\item
  Plumillas de longitud 6 cm y 6 cm respectivamente
\item
  Ángulo de apertura 90°
\end{itemize}

Este es el dibujo que corresponde exactamente a las medidas descritas en
el problema inicial.

\subsection{Answerlist}\label{answerlist-1}

\begin{itemize}
\tightlist
\item
  ``Verdadero''
\item
  ``Falso''
\item
  ``Falso''
\item
  ``Falso''
\end{itemize}

\section{Meta-information}\label{meta-information}

exname: parabrisas\_dimensiones\_esquema\_v1 extype: schoice exsolution:
1000 exshuffle: TRUE exsection: Geometría y representación gráfica

\end{document}
